\SourceLeaf{402}

Um die Derivirte nach $y$ von jeder einzelnen der vier Grössen $\alpha,\allowbreak \beta,\allowbreak \gamma,\allowbreak \delta$ zu bestimmen, wird zunächst erforderlich sein, jenen Ausdruck so zu verwandeln, dass er nur zwei der Grössen enthält. Dies kann z.\ B. mit Hülfe der im vorhergehenden Artikel bestimmten nach $\eta$ genommenen zweiten Derivirten der Quotienten geschehen, so dass also alle Ausdrücke
\[
\begin{aligned}&\dfrac{\mathrm d\log\dfrac{\alpha}{\beta}}{\mathrm d\log y}
-\dfrac{\mathrm d\mathrm d\log\dfrac{\alpha}{\beta}}{(\mathrm d\log\eta)^2}
\\[8pt]&{}-\dfrac12\dfrac{\mathrm d\log\dfrac{\alpha\beta}{xx}}{\mathrm d\log\eta}
\cdot\dfrac{\mathrm d\log\dfrac{\alpha}{\beta}}{\mathrm d\log\eta}\end{aligned}
\]
verschwinden, in welcher Reihenfolge $\alpha,\allowbreak \beta,\allowbreak \gamma,\allowbreak \delta$ mit einander auch vertauscht werden und welcher gemeinsame Index ihnen und den Zeichen $y,\allowbreak x,\allowbreak \eta$ gegeben werde. Ersetzt man $\frac{\alpha}{\beta}$ durch $\frac{b'\beta'}{\beta\beta}$, $\alpha\beta$ durch $b'\beta'$ und nach dem vorhergehenden Artikel
\[
\begin{gathered}\dfrac{\mathrm d\log\dfrac{b'}{k}}{2\mathrm d\log y}
\quad\text{durch}\\[9pt]
\dfrac{\mathrm d\mathrm d\log\dfrac{\beta'}{x}}{(2\mathrm d\log\eta)^2}
+\dfrac12\left[\dfrac{\mathrm d\log\dfrac{\beta\beta}{\beta'}}{\mathrm d\log\eta}\right]^2\end{gathered}
\]
so ersieht man, dass der Ausdruck
\[
\begin{aligned}&\dfrac{\mathrm d\log\dfrac{\beta}{x}}{\mathrm d\log y}
-\dfrac{\mathrm d\mathrm d\log\dfrac{\beta}{x}}{(\mathrm d\log\eta)^2}
\\[8pt]&{}-\dfrac12\left[\dfrac{\mathrm d\log\dfrac{\beta}{x}}{\mathrm d\log\eta}\right]^2\end{aligned}
\]
seinem Werthe nach ungeändert bleibt, nicht nur, wie eben gefunden, wenn $\beta$ mit $\alpha$ oder $\gamma$ oder $\delta$ vertauscht wird, sondern auch wenn diesen Grössen zugleich mit $y,\allowbreak x,\allowbreak \eta$ ein beliebiger gemeinsamer Index gegeben wird. Für einen beständig wachsenden Index verschwindet der Werth des Ausdrucks und durch Multiplication mit $\frac12\surd {\frac{\beta}{x}}$ entsteht demnach
\[
\dfrac{\mathrm d\surd {\dfrac{\alpha}{x}}}{\mathrm d\log y}
-\dfrac{\mathrm d\mathrm d\surd {\dfrac{\alpha}{x}}}{(\mathrm d\log\eta)^2}=0
\]
als partielle Differentialgleichung für jede der Grössen $\alpha,\allowbreak \beta,\allowbreak \gamma,\allowbreak \delta$.]

\AGMArticle{22}
[Alle Glieder des Algorithmus sind ihrem Werthe nach abhängig von vier Grössen, als solche dürfen $a,\allowbreak b,\allowbreak \alpha,\allowbreak \beta$ angenommen werden, aber auch $k,\allowbreak y,\allowbreak x,\allowbreak \eta$ weil diese unabhängig von einander sich ändern können. Sämmtliche Gleichungen zwischen den Gliedern der Reihen lassen sich in solche Form bringen, dass sie sowol in Bezug auf $a,\allowbreak b,\allowbreak c,\allowbreak .\;. a^{\mathrm n},\allowbreak b^{\mathrm n},\allowbreak c^{\mathrm n},\allowbreak k$ als auch in Bezug auf


\SourceLeaf{403}

\[
\begin{gathered}\dfrac{\alpha}{k},\ \dfrac{\beta}{k},\ \dfrac{\gamma}{k},\ \dfrac{\delta}{k}\ .\;.\;.\\[8pt]\
{}^{2^n}\!\surd {\dfrac{\alpha^{\mathrm n}}{k}},\
{}^{2^n}\!\surd {\dfrac{\beta^{\mathrm n}}{k}},\
{}^{2^n}\!\surd {\dfrac{\gamma^{\mathrm n}}{k}},\
{}^{2^n}\!\surd {\dfrac{\delta^{\mathrm n}}{k}},\ \dfrac{x}{k}\end{gathered}
\]
homogen sind. Die Verhältnisse zwischen diesen letzteren hangen also nicht von $k$ und $x$, sondern allein von $y$ und $\eta$ ab, wir können also
\[
\begin{gathered}\alpha=x P(y,\eta)^2,\\[5pt]
\beta=x Q(y,\eta)^2,\\[5pt]
\gamma=x R(y,\eta)^2,\\[5pt]
\delta=x S(y,\eta)^2\end{gathered}
\]
setzen und aus den Betrachtungen in Art.\ 18 folgt dann, dass diese Functionen $P,\allowbreak \ Q,\allowbreak \ R,\allowbreak \ S$ dieselben bleiben, wenn $\alpha,\allowbreak \beta,\allowbreak \gamma,\allowbreak \delta,\allowbreak x,\allowbreak y,\allowbreak \eta$ mit einem beliebigen Index versehen werden. Aus demselben Artikel folgt auch, mit Benutzung der in Artikel 15 gebrauchten Bezeichnung, für die mit einem gleichen Index behafteten $a,\allowbreak b. c,\allowbreak y$
\[
\begin{aligned}
P(y,1)&=py=\surd {\dfrac ak}=\surd {\dfrac a{\mathrm M(a,b)}}\\
Q(y,1)&=qy=\surd {\dfrac bk}\\
R(y,1)&=ry=\surd {\dfrac ck}\\
S(y,1)&=0
\end{aligned}
\]
ferner als Gleichungen, welche denjenigen entsprechen, durch die die Grössen $\gamma$ und $\delta$ bestimmt wurden:
\[
\begin{gathered}
qqrr(qqRR-rrQQ)\\[4pt]=rrpp(ppRR-rrPP)\\
=ppqq(qqPP-ppQQ)\\[4pt]=ppqqrrSS\\[4pt]
\dfrac{PP+QQ}{p(yy)}=\dfrac{RR-SS}{r(yy)}.\\[9pt]
\dfrac{PP-QQ}{r(yy)}=\dfrac{RR+SS}{p(yy)}\\[4pt]
\dfrac{PP+RR}{p(\surd  y)}=\dfrac{QQ+SS}{q(\surd  y)},\\[9pt]
\dfrac{PP-RR}{q(\surd  y)}=\dfrac{QQ-SS}{p(\surd  y)}
\end{gathered}
\]
und als Gleichungen, die das Bildungsgesetz des Algorithmus darstellen:
\[
\begin{aligned}
2p(yy)\cdot P(yy,\eta\eta)&=PP+QQ,\\[5pt]
2p(yy)\cdot R(yy,\eta\eta)&=RR+SS\\[5pt]
2r(yy)\cdot R(yy,\eta\eta)&=PP-QQ,\\[5pt]
2r(yy)\cdot P(yy,\eta\eta)&=RR-SS\\[5pt]
q(yy)\cdot Q(yy,\eta\eta)&=PQ,\\[5pt]
q(yy)\cdot S(yy,\eta\eta)&=RS
\end{aligned}
\]
worin alle Functionen, neben denen kein Argument geschrieben ist, sich auf die einfachen $y$ und $\eta$ beziehen.

\SourceLeaf{404}

Zur Berechnung der Werthe der Functionen, welche gegebenen Werthen von $\frac ab=\frac{pp}{qq}$ und $U$ oder $V$ zugehören, erhält man aus der \textsc{Gauss}ischen Formel für $H$ oder $x$, wenn man zur Abkürzung
\[
\begin{aligned}&(a^{\mathrm n}\sin 2^{\mathrm n} U^{\mathrm n})^2+(b^{\mathrm n}\cos 2^{\mathrm n} U^{\mathrm n})^2\\[5pt]&=b^{\mathrm n} b^{\mathrm n}\Delta^{\mathrm n}\Delta^{\mathrm n}\end{aligned}
\]
setzt, die Gleichungen:
\[
\begin{aligned}
Q&=q\cdot\surd \Delta\cdot{}^{4}\!\surd {\Delta'}\cdot{}^{8}\!\surd {\Delta''}\cdot{}^{16}\!\surd {\Delta'''}.\;.\;.\\[9pt]
P&=\dfrac pq Q\dfrac1\Delta\\[9pt]
R&=\dfrac rq Q\dfrac{\cos U}{\Delta}\\[9pt]
S&=\dfrac pr Q\dfrac{\sin U}{\Delta}i
\end{aligned}
\]

Die Bestimmung der Grössen $k,\allowbreak y,\allowbreak x,\allowbreak \eta$ als Grenzwerthe lässt erkennen, dass, bei geeigneter hier noch zulässiger Wahl der Vorzeichen der Functionen $P,\allowbreak Q,\allowbreak R,\allowbreak S$, für bis zur Null abnehmende Werthe von $y,\allowbreak y\eta,\allowbreak \frac1\eta$ die Ausdrücke
\[
\begin{gathered}P(y,\eta),\\[6pt] Q(y,\eta).\\[6pt]
\dfrac{R(y,\eta)}{y^{\frac14}\eta^{\frac12}},\\[6pt]
\dfrac{S(y,\eta)}{y^{\frac14}\eta^{\frac12}}\end{gathered}
\]
sich dem Grenzwerthe Eins nähern, dass aber für ein beständig abnehmendes $y$ und ein endliches reelles $u$ die Ausdrücke
\[
\begin{gathered}P(y,e^{2iu}),\\[6pt] Q(y,e^{2iu}),\\[6pt]
\dfrac{R(y,e^{2iu})}{2y^{\frac14}\cos u},\\[6pt]
\dfrac{S(y,e^{2iu})}{i2y^{\frac14}\sin u}\end{gathered}
\]
jenen Grenzwerth haben.

Die Functionen $P,\allowbreak Q,\allowbreak R,\allowbreak iS$ haben reelle Werthe für ein complexes $\eta^{\frac12}$ von der Form $e^{iu}$ und bleiben bis auf $iS$, welches nur sein Zeichen wechselt, ungeändert, wenn man $u$ in $-u$ verwandelt, es ist also
\[
\begin{aligned}\dfrac{P(y,\dfrac1\eta)}P
&=\dfrac{Q(y,\dfrac1\eta)}Q
\\[7pt]&=\dfrac{R(y,\dfrac1\eta)}R
\\[7pt]&=\dfrac{-S(y,\dfrac1\eta)}S=1
\text{]}\end{aligned}
\]

\AGMArticle{23}
[Bildet man denselben Algorithmus wie vorher für $\alpha,\allowbreak \beta,\allowbreak \gamma,\allowbreak \delta$, jetzt für $A,\allowbreak B,\allowbreak C,\allowbreak D$ und macht $A=\gamma,\allowbreak \ B=\delta$, so wird offenbar für jedes $n$, $A^{\mathrm n}=\gamma^{\mathrm n},\allowbreak \ B^{\mathrm n}=\delta^{\mathrm n},\allowbreak \ C^{\mathrm n}=\alpha^{\mathrm n},\allowbreak \ D^{\mathrm n}=\beta^{\mathrm n}$ und also für die Grenzwerthe $K,\allowbreak H$, welche den
